\paragraph{Small representatives of the three rational spans.}\label{par:shifted-small}
Write \(L=x^\epsilon\), \(z=frL\), \(h=frn\), and \(d\) for the
total degree of the harmonic polynomial at weight \(u\). Let \(A\)
be the sum of all the flag entries of \(u\), so that \(d\le A\) and
\(A+h+1\le |u|\). We use the selected size bounds
\[
 L\ge2,\qquad \log f,\ \log\operatorname{rd}(E),\
            \log(1+A)\le2L,\qquad \log n\le L,\qquad r\le n .
\]
They hold at each of the three weights separately by coordinatewise
domination by \(\kappa\). We give estimates independent of the
lattice class and the embedding of \(K\).

\begin{proposition}[Individual height bounds at the three weights]
\label{prop:three-shifted-spanning-sets}
Each of the three specified weights has a rational spanning set
in its scaled theta lattice in which each chosen form \(g\) satisfies
\[
 H(s_uM_0^{frn}g)\le G_c|u|.
\]
Here \(H\) uses the normalization of the proof. The size bound is independent of the
source lattice class and the coefficient embedding and is retained by
any subset chosen as a basis.
\end{proposition}

\begin{proof}
On a unimodular source lattice \(\Lambda\) of rank \(n\), choose
\(2fn\) rationally independent vectors whose norm at every embedding
is at most \(\rho=2f^{3/2}n^2\operatorname{rd}(E)\).
Indeed for \(v\ne0\) in its dual for the trace of the Hermitian
pairing, \(\langle v,\Lambda\rangle\subseteq\mathfrak D_E^{-1}\).
Unimodularity then implies \(\mathfrak D_E v\subseteq\Lambda\).
Hence \(\langle v,v\rangle\in\mathfrak D_E^{-2}\), using invariance
of the different under conjugation. Its positive conjugates have
product at least \(|\operatorname{disc}E|^{-2}\). Their sum is
at least \(2f/\operatorname{rd}(E)^2\).

For the \(2f\) positive conjugates \(a_\sigma\), the arithmetic--geometric mean bound gives
\[
  \sum_\sigma a_\sigma\ge 2f\Big(\prod_\sigma a_\sigma\Big)^{1/(2f)}\ge \frac{2f}{\operatorname{rd}(E)^2}.
\]
 Euclidean transference in
dimension \(2fn\) bounds the last successive minimum of the
original lattice by \(2fn\,\operatorname{rd}(E)/\sqrt{2f}\)
in the trace norm, which suffices for the displayed \(\rho\).
The source lattices are full and finitely generated since the
admissible-level conditions sandwich them between two scalar
multiples of the standard lattice.
\cite{banaszczyk}\label{cor:slice-short-heights-1}

\proofstep{The bounded harmonic grid.}
Use these vectors for the two sets of row letters in the harmonic
polynomial of the short rule. Here is a finite interpolation argument
which controls their size even for very large row weights. Write
\(a+b=r\) for the two block sizes. Start with three lists of lengths
\(a,a,b\) of vectors \(e_i,e'_i,g_j\). Let 
\[
  B_{li}=\langle
e'_l,e_i\rangle,
\]
 with the pairing linear in the first variable.

In the second block use the letters
\[
 f_j=(\det B)g_j-\sum_{i,l}\langle g_j,e_i\rangle
                         (\operatorname{adj}B)_{il}e'_l .
\]
They are integral and \(\langle f_j,e_i\rangle=0\) identically,
including for singular \(B\). Use the products of leading minors
with the row-gap exponents, on \(e\) and the conjugates of \(f\).
They satisfy the flag weight and are harmonic because each mixed
contraction pairs orthogonal letters. In the opposite convention
conjugate both the lattice and its letters before making the
construction. This gives the integral polynomial in the symmetric
algebra of the original lattice.

\proofstep{Spanning after extension to the splitting field.}
For spanning, over a splitting field the independent embeddings
and their conjugates are the separate factors of

\[
  E\otimes_{\Q}\overline{\Q}.
\]
 The simultaneous orbit of a highest
minor product spans its highest-weight harmonic space under column
changes: on each factor it generates the irreducible trace-free
Schur representation with the two prescribed partitions. The
condition \(a+b\le n\) permits the two opposite blocks of an
invertible column frame and its dual. Such orthogonal pairs on the
open orbit are obtained in the displayed parametrization by choosing
the primed list dual to the first block and setting the last list
equal to the desired perpendicular vectors. This prescription can
be made in the split factors because all the signed embeddings
are distinct.

Express the three lists in the rational basis of
\(2fn\) short vectors. The map from coefficients to the minor
product is polynomial of degree at most 
\[
  D=(2r+1)d:
\]
the adjugate expression has degree at most \(2a+1\). Any linear
functional annihilating its values at coordinates in
\(\{0,\ldots,D\}\) annihilates the polynomial, by successive
univariate interpolation. The grid therefore spans the same
space in characteristic zero. Degree zero uses the constant
polynomial when present, and the zero space uses the empty set.

With \(d>0\), each grid vector has length at most

\[
  R=2fnD\rho
\]
 at every embedding. Cauchy--Schwarz and determinant
expansion bound the length of every projected letter by
\((a+1)!R^{2a+1}\). Each linear form on \(n\) coordinates has
coefficient sum at most \(\sqrt n\) times its Euclidean norm.
The determinant expansion in a leading minor of size \(i\) costs
at most \(i!\le r^i\). Submultiplicativity thus gives, at every
embedding,
\[
 \|P\|_1\le\{r\sqrt n\,(r+1)!\,R^{2r+1}\}^{d},
 \qquad \log\|P\|_1\le140rLd .
\]
Indeed \(R\le e^{17L}\), \((r+1)!\le e^{4rL}\), and
\(r\sqrt n\le e^{2L}\) under the size bounds above. The rule is
linear in \(P\) for fixed lattice, marking and translates by the
covariant derivative expression on each chart. Replacing its
polynomial by this grid family preserves the rational span.
\label{cor:slice-short-heights-2}

\proofstep{The passage from coefficient height to Petersson norm.}
Here are details of the passage from coefficient size to the
supremum at an embedding. For a nonzero section choose an
orthonormal frame at more than half the supremum of its invariant
norm before the Lie algebra normalization. At that frame the first
minimum \(\lambda\) of the period lattice satisfies
\(\log(1/\lambda^2)\le80L\).
In the rank-one period argument of the integral lower norm, take
spread bound \(w\), transvection period \(M\), reference-lattice
clearing integer \(N_A\), and polarization factor \(k_0\le2\).
The period bound is
\[
 \lambda^2\ \ge\ k_0 e_0/N_A^2,\qquad
 e_0=\min\{1,\pi/[M(fw+2)k_0\log 2]\}.
\]

In fact \(N_A v\) belongs to the reference lattice and if every
majorant of that vector were below \(e_0\), the negative imaginary
shift by \(1/(2k_0e_0)\) and Schwarz's lemma at the vanishing cusp
would contradict the half-supremum choice. In the definite case a
nonzero isotropic vector is already excluded. The rank-one data of
the run give \[
 \log M\le30L,\ \log(fw+2)\le5L,\ \log N_A\le4L,\qquad \log(1/k_0)\le3L.
\] Thus \(-\log e_0\le36L\), which gives
the asserted bound. Conjugate embeddings use the conjugate
uniformization with the same norm.
\label{cor:slice-short-heights-3}

In normalized Gaussian coordinates the tensor period lattice has
minimum at least \(\lambda/\rho\). A nonzero tensor has nonzero
contraction with one of the \(2fn\) short vectors, and that contraction
is a period of length at most \(\rho\) times the tensor norm. Its
torus is principally polarized and has covolume one in these
coordinates. For a theta function of mean squared norm one in
dimension \(h\), the polydisc of radius
\[
 b_0=\min\{1/\sqrt{\pi h},\lambda/(4\rho\sqrt h)\}
\]
in every complex coordinate injects into the torus. Translate the
function by the unitary Heisenberg operator to put its evaluation
at zero. The submean bound for its squared modulus uses the
polydisc volume \((\pi b_0^2)^h\). The Gaussian density changes
on that polydisc by at most \(e\).

Taking square roots in the submean estimate, the invariant supremum of the normalized function \(h_0\) satisfies
\[
  \|h_0\|_{\sup}\le\left(\frac{e}{(\pi b_0^2)^h}\right)^{1/2}=e^{1/2}(\sqrt\pi\,b_0)^{-h}.
\]
The logarithm of the last factor is at most \(55hL\), by the
bounds on \(\lambda,\rho,f,r,n\). For a derivative of multi-index
\(\alpha\) of degree \(d\), Cauchy's formula in the active variables
with radii \(\sqrt{\alpha_i/\pi}\) costs at most
\((\pi e d)^{d/2}\) times that supremum, since
\(\alpha_i!\le\alpha_i^{\alpha_i}\). A homogeneous polynomial costs
its coefficient sum times this factor. Changing from normalized
letters to the original letters multiplies this by \(t^{d/2}\),
where \(t=2\xi_R\) is the polarization scale. The Lie norm of the
run contributes \((\pi t)^{-d/2}\).
The factors involving the polarization scale \(t\) cancel, leaving
\[
  t^{d/2}(\pi t)^{-d/2}=\pi^{-d/2}.
\]

\proofstep{The unit-mean normalization.}
For clarity, the unit mean normalization has no class-index cost in
the short reading. Write \(J=[\Lambda:\Lambda\cap\mathcal O_E^n]^r\)
with the index taken on the side specified by the tensor translate
count. Equivalently it is the number of cosets in the projected
vacuum sum before averaging. Distinct cosets have orthogonal
translate characters on a common finite cover, so their sum has
squared mean \(J\) times that of the tensor power of the initial
vacuum. Projection divides the sum by \(J\). After removal of the
initial vacuum metric its norm consequently has the extra factor
\(1/\sqrt J\). In the rule the torsion count divided by its initial
level index equals \(J\), whereas the arithmetic normalizing
factor \(\beta_\Lambda\) has \(|\sigma\beta_\Lambda|^2=J\) and
occurs inversely. The resulting factor is

\[
  J/(\sqrt J\sqrt J)=1.
\]
 The residual finite-character scalar
has modulus one. The finite mean of translates can be bounded
term by term, since all its coefficients have modulus at most
\(p_{\rm type}^{\,h}\), with \(\log p_{\rm type}\le L\).
These are the two normalizations in the short rule on the
marked level.
\label{cor:slice-short-heights-4}

Including the factor two from evaluation near the supremum, the
preceding estimates give each unscaled short representative the
bound
\[
 \log\|g\|_{\sup,\sigma}
 \le 140rLA+10LA+200hL\le400z|u|.
\]
In particular no count of the potentially large interpolation
grid occurs in this bound. Finite dimension selects a finite spanning
set \(\operatorname{Sh}(u)\) of the separate short representatives.

\proofstep{Clearing the product lattice.}
To put these forms in the lattice used for products, let
\(M_0=g_0g_1\), and let \(s_u\) be the selected nonzero weight
scalar, namely \(s_F,s_0,s_\kappa\) at the respective weights.
The short rule contributes \(M_0^{frn}g\) to the unscaled
theta lattice used here. Thus the actual integral columns
are \(s_uM_0^{frn}g\). Their rational span is unchanged.

The common scale bounds \(\log M_0\le\lambda_*\) and
\(\log|\sigma s_u|\le2(\lambda_*+2\log p)|u|\), where \(p\)
is the run prime, give the mean log bound
\[
 H(s_uM_0^{frn}g)
 \le [400z+2(\lambda_*+2\log p)+\lambda_*]\,|u|
 \le G_c |u|.
\]
For the last inequality the growth allowance contains \(1009z\)
and \(2fr(\lambda_*+2\log p)\), as well as the nonnegative
coefficient-field and lattice terms. Here \(fr\ge 12\).
Every subset selected by independence at a later stage inherits
the same individual bound.
\label{cor:slice-short-heights-5}
\end{proof}
